在前几部分中，我们主要从技术角度考察智能体大语言模型，讨论其架构、训练过程、推理策略、工具使用及不同领域的应用。本部分转换视角，转而探究智能体大语言模型在社会情境中如何\emph{表现}。
随着智能体大语言模型越来越多地以助手、协作者、导师等身份进入社会情境，它们的行为倾向开始具有实际意义。例如，智能体在策略互动中倾向于合作还是背叛，是否会对礼貌作出回应，以及能否在多智能体模拟中追踪其他智能体的信念状态，都会产生现实影响。

本部分讨论的大多数研究借鉴了最初用于研究人类及其他动物的方法，考察策略行为、社交礼仪、共情、人格等特征。每章末尾都简要反思一个共同问题：将这些方法用于人工智能，能否取得有意义的结果？

在下一部分，也就是最后一部分专题内容中，我们将进一步探讨大语言模型的内部机制，以及各种认知形式乃至意识是否可能涌现。
